% Claim proof file for row claim_3_8 -- "DisjointHardInterventions".
%
% Proof lifted from lecture-notes/lecture_notes/graphs.tex:505--533 (the
% \Claude{...} block immediately following the claim). Only light edits:
% explicit \refrow{def_3_10} and \refrow{def_3_11} cross-links the first
% time \doit and \spl appear in the body, \medskip\noindent\emph{...}
% section headers mirroring claim_3_4_proof_HardInterventionsCommute.tex,
% a \checkmark at the end of each component-agree paragraph, and a closing
% one-liner. The \Claude{...} wrapper from the LN source is dropped.
%
% Compiles standalone via the `subfiles` package.

\documentclass[main]{subfiles}

\begin{document}
% ref: claim_3_8 (proof, with statement restated for readability)
% title: DisjointHardInterventions
\def\rowref{claim\_3\_8}
\phantomsection\label{claim_3_8}
%
% Cross-references: see claim_<ref>_statement_<title>.tex in this folder
% for the corresponding statement.

% --- statement (restated) ---------------------------------------------------
\begin{Lem}[Disjoint hard interventions and node-splittings commute]
   Let $G=(J,V,E,L)$ be a CDMG and $W_1 \ins J \cup V$ and $W_2 \ins V$ two disjoint subsets of nodes of $G$.
   Then the CDMG obtained from first hard intervening on $W_1$ and then node-splitting $W_2$ is the same CDMG that arises from first node-splitting $W_2$ and then hard intervening on $W_1$.
      \[ \lp G_{\doit(W_1)} \rp_{\spl(W_2)} =  \lp G_{\spl(W_2)} \rp_{\doit(W_1)}.   \]
\end{Lem}

% --- proof ------------------------------------------------------------------
\begin{proof}
\medskip\noindent\emph{Preamble.}
Since $W_1 \ins J \cup V$ and $W_2 \ins V$ are disjoint, both operations are
well-defined in either order: on the left, $\spl(W_2)$ acts on $G_{\doit(W_1)}$
whose output set $V \sm W_1$ still contains $W_2$ (using $W_1 \cap W_2 = \emptyset$);
on the right, $\doit(W_1)$ acts on $G_{\spl(W_2)}$, whose input/output universe
contains every vertex of $W_1 \ins J \cup V$ untouched by the split (again using
$W_1 \cap W_2 = \emptyset$). For $w \in W_2$ we write $w^0$ (incoming copy) and
$w^1$ (outgoing copy) for the two split copies of $w$; for $v \notin W_2$ we
identify $v^0 = v^1 = v$. Each side of the displayed equation is a CDMG presented
as a 4-tuple $(J_\ast, V_\ast, E_\ast, L_\ast)$, and we verify the four components
in turn.

\medskip\noindent\emph{Node sets.}
Left: by \refrow{def_3_10}, $G_{\doit(W_1)}$ has input set $J \cup W_1$ and
output set $V \sm W_1$; then $\spl(W_2)$~\refrow{def_3_11} (with
$W_2 \ins V \sm W_1$, by $W_1 \cap W_2 = \emptyset$) gives
\[
  J \cup W_1
  \qquad\text{and}\qquad
  (V \sm W_1 \sm W_2) \dcup W_2^0 \dcup W_2^1.
\]
Right: $G_{\spl(W_2)}$ has input set $J$ and output set $(V \sm W_2) \dcup W_2^0 \dcup W_2^1$;
then $\doit(W_1)$ (with $W_1 \cap W_2^0 = W_1 \cap W_2^1 = \emptyset$, since
$W_2^0, W_2^1$ are fresh split copies of vertices in $W_2$ and $W_1 \cap W_2 = \emptyset$)
gives
\[
  J \cup W_1
  \qquad\text{and}\qquad
  (V \sm W_1 \sm W_2) \dcup W_2^0 \dcup W_2^1.
\]
These agree. \checkmark

\medskip\noindent\emph{Directed edges.}
Left: $E_{\doit(W_1)} = E \sm \lC v \tuh w \mid w \in W_1 \rC$; then $\spl(W_2)$
yields
\[
  \lC v_1^1 \tuh v_2^0 \mid v_1 \tuh v_2 \in E,\ v_2 \notin W_1 \rC
  \cup \lC w^0 \tuh w^1 \mid w \in W_2 \rC.
\]
Right: $E_{\spl(W_2)} = \lC v_1^1 \tuh v_2^0 \mid v_1 \tuh v_2 \in E \rC \cup \lC w^0 \tuh w^1 \mid w \in W_2 \rC$;
then $\doit(W_1)$ removes the edges with head in $W_1$. An edge $v_1^1 \tuh v_2^0$
has head $v_2^0 \in W_1$ iff $v_2 \in W_1$ (using $W_1 \cap W_2 = \emptyset$: if
$v_2 \in W_2$ then $v_2^0$ is a fresh split copy not in $W_1$, and if $v_2 \notin W_2$
then $v_2^0 = v_2$). The split edges $w^0 \tuh w^1$ have head $w^1 \notin W_1$,
since $w \in W_2$ and $W_1, W_2$ are disjoint. After removal we obtain
\[
  \lC v_1^1 \tuh v_2^0 \mid v_1 \tuh v_2 \in E,\ v_2 \notin W_1 \rC
  \cup \lC w^0 \tuh w^1 \mid w \in W_2 \rC,
\]
which agrees with the left-hand side. \checkmark

\medskip\noindent\emph{Bidirected edges.}
Left: $L_{\doit(W_1)} = L \sm \lC v \huh w \mid w \in W_1 \rC$; then $\spl(W_2)$
yields
\[
  \lC v_1^0 \huh v_2^0 \mid v_1 \huh v_2 \in L,\ v_1, v_2 \notin W_1 \rC.
\]
Right: $L_{\spl(W_2)} = \lC v_1^0 \huh v_2^0 \mid v_1 \huh v_2 \in L \rC$;
then $\doit(W_1)$ removes the bidirected edges touching $W_1$. Using
$v_k^0 \in W_1 \iff v_k \in W_1$ for $k = 1, 2$ (same fresh-copy / identity case
analysis as for the directed edges, again using $W_1 \cap W_2 = \emptyset$), the
result is
\[
  \lC v_1^0 \huh v_2^0 \mid v_1 \huh v_2 \in L,\ v_1, v_2 \notin W_1 \rC,
\]
which agrees with the left-hand side. \checkmark

All four components agree, completing the proof.
\end{proof}

\end{document}
